\documentclass[10pt,a4paper]{amsart}
\usepackage[margin=19mm]{geometry}
\usepackage{amsmath,amssymb,mathtools}
\usepackage[scr=rsfs,cal=euler]{mathalfa}
\usepackage{xfrac}
\usepackage[unicode,hidelinks,bookmarks=false]{hyperref}
\newcommand{\ie}{i.e.\ }
\newcommand{\cf}{cf.\ }
\newcommand{\resp}{respectively\ }
\newcommand{\Subsection}{\subsection}
\providecommand{\nobreakdash}{\nobreak\hbox{-}}
\setlength{\emergencystretch}{2em}
\multlinegap0pt
\usepackage{bm}
\usepackage{bbm}
\usepackage{dsfont}
\usepackage{xspace}
\usepackage{cancel}
\usepackage{mathtools}
\usepackage[shortlabels]{enumitem}
\usepackage{amsthm}
\makeatletter
\@ifundefined{corollary}{\newtheorem{corollary}{Corollary}}{}
\@ifundefined{lemma}{\newtheorem{lemma}{Lemma}}{}
\@ifundefined{proposition}{\newtheorem{proposition}{Proposition}}{}
\@ifundefined{remark}{\newtheorem{remark}{Remark}}{}
\@ifundefined{theorem}{\newtheorem{theorem}{Theorem}}{}
\makeatother
\newcommand{\cal}[1]{\mathcal{#1}}
\title[On Lattice Points, Short-Time Estimates, and Global Well-pos]{On Lattice Points, Short-Time Estimates, and Global Well-posedness of the Quintic NLS on $\mathbb{T}$}
\author{Ryan McConnell}
\date{Source: arXiv:2310.01638v4 (CC BY 4.0)}
\begin{document}
\maketitle

\noindent\emph{Source and licence.} On Lattice Points, Short-Time Estimates, and Global Well-posedness of the Quintic NLS on $\mathbb{T}$. Version arXiv:2310.01638v4 (CC BY 4.0), distributed under the Creative Commons Attribution 4.0 International licence, as recorded in the arXiv licence metadata for this exact version and shown on the abstract page https://arxiv.org/abs/2310.01638v4. Full rights notice: licenses/arxiv-2310.01638-CC-BY-4.0.txt.\par\smallskip

\noindent\textit{Editorial scope.} Counted unit: the Lemma whose source label is \texttt{Lemma: M6 integral term} in the article (source0.tex lines 1442--1447, source label \texttt{Lemma: M6 integral term}), delivered together with the article's own complete original proof of it (source0.tex lines 1448--1619, one \texttt{proof} environment of 172 lines). No mathematics is written, repaired or re-derived: statement and proof are the article's own text line for line. Required context retained uncounted: Equation: Quintic NLS (lines 45--50); Theorem: GWP theorem (lines 94--96); Proposition: Trilinear Free Estimate (lines 702--719); Condition: M condition (lines 712--714); Lemma: trilinear Estimates (lines 886--904); Corollary: Trilinear Strichartz Specific (lines 950--960); Equation: Regularity (lines 1101--1103); Lemma: M6 estimate (lines 1122--1147); equation: M bound 1 (lines 1130--1132); Equation: M bound 2 (lines 1135--1137); Equation: M bound 3 (lines 1139--1141); Equation: M bound 4 (lines 1143--1145); Remark: N12 M term (lines 1179--1185). Every definition, display and label of those lines is retained unchanged. Omitted from this package and not redistributed: 1--44, 51--93, 97--701, 715--885, 905--949, 961--1100, 1104--1121, 1133--1134, 1138--1138, 1142--1142, 1146--1178, 1186--1441, 1620--1798. Nothing inside a retained excerpt is omitted or altered: every retained body is delivered exactly as the article's lines are written, so no omission receipt of any kind accompanies this package. Citations: the retained excerpts carry no citation command at all, so no bibliography entry of the article is transcribed and no reference is redistributed. Reference adaptation: none was needed and none was made; every reference inside the delivered unit resolves inside it, so no unresolved reference or citation is produced. Printed slips kept literal: no printed word, symbol, hypothesis or number of the counted statement or of its proof was altered. Layout and class: the article's own class is \texttt{amsart} with packages including fullpage, tikz, amsmath, hyperref, tcolorbox, amsfonts, amssymb, amsthm, enumitem, fancyhdr, color, soul, comment, standalone; the wrapper is the lane's pinned \texttt{amsart} editorial wrapper at 10pt on A4, which restates the theorem-like environments the retained text needs and adds the article's own macro definitions that the retained text needs (\texttt{\textbackslash cal}), with extra wrapper packages (bm, bbm, dsfont, xspace, cancel, mathtools, enumitem). The delivered numbering is the unit's own. Assets: no retained span contains a figure environment, a graphics-inclusion command, a PDF-page inclusion, a table or a TikZ picture; no image file is copied into this package, the package declares no asset dependency and no image file is redistributed with it. Licence: version arXiv:2310.01638v4 of this article is distributed under the Creative Commons Attribution 4.0 International licence, as recorded in the arXiv licence metadata for this exact version and shown on the abstract page https://arxiv.org/abs/2310.01638v4. A full rights notice is delivered at licenses/arxiv-2310.01638-CC-BY-4.0.txt. Characters: every retained excerpt is pure ASCII, so the delivered engine, XeLaTeX with its default Latin Modern text font, reports no missing character.
\begin{equation}\label{Equation: Quintic NLS}
    \begin{cases}
        (i\partial_t+\Delta)u = \pm |u|^4u\\
        u(x,0) = u_0(x)\in H^s(\mathbb{T}),
    \end{cases}
\end{equation}

\begin{theorem}\label{Theorem: GWP theorem}
    Let $\alpha > \frac{131}{208}$ and $s > \frac{\alpha}{3}$. Then \eqref{Equation: Quintic NLS} is globally well-posed in $H^s(\mathbb{T})$ for any $\|u_0\|_{L^2_{x}(\mathbb{T})}\ll 1$.
\end{theorem}

\begin{proposition}\label{Proposition: Trilinear Free Estimate}
    Let $\lambda\gg 1$, $I_1, I_2, I_3\subset \frac{1}{\lambda}\mathbb{Z}$ be dyadic intervals such that $|I_1|\lesssim |I_2|\lesssim |I_3|$. If $N_{\max}$ is defined by
    \[
    N_{\max} = \max(N_{13}, N_{23})\gg 1, \quad N_{ij} = \mathrm{dist}(I_i,I_j),\quad 1\leq i\leq  j\leq 3,
    \]
    and satisfies $N_{\max}\gg |I_2|$, then for any $\phi_j\in L^2(\mathbb{T})$, $1\leq j\leq 3$, with $\mathrm{supp}\,\widehat{\phi}_j\subset I_j$ we have the trilinear estimate:
    \begin{equation*}\label{Equation: Trilinear base}
\|\prod_{j=1}^3e^{it\Delta}\phi_j\|_{L^2_{x,t}(\lambda\mathbb{T}\times\mathbb{T})}\lesssim \left(\frac{1}{\lambda} + \frac{|I_2|}{N_{\max}}\right)^\frac{1}{2}\prod_{j=1}^3\|\phi_j\|_{L^2_x(\lambda\mathbb{T})}.
    \end{equation*}
    Alternatively, if $N_{\max}\gtrsim |I_2|$ and $|I_1|\ll|I_2|$, $I_1-I_3\subset B(0, J)$ for some $J > 0$, and
 \begin{equation}\label{Condition: M condition}
M := \frac{|I_1|(J+|I_1|)}{N_{23}}\ll |I_2|\wedge N_{23},
 \end{equation}
    then we find the enhanced trilinear estimate:
    \begin{equation*}\label{Equation: Trilinear HHL}
\|\prod_{j=1}^3e^{it\Delta}\phi_j\|_{L^2_{x,t}(\lambda\mathbb{T}\times\mathbb{T})}\lesssim \left(\frac{1}{\lambda} + \frac{M\vee |I_1|}{N_{\max}}\right)^\frac{1}{2}\prod_{j=1}^3\|\phi_j\|_{L^2_x(\lambda\mathbb{T})}.
    \end{equation*}
\end{proposition}

\begin{proposition}\label{Lemma: trilinear Estimates}
    Suppose $s < \frac{\alpha}{2}$, and let $L\in 2^{\mathbb{N}_0}$, $I_1, I_2, I_3$ be intervals contained in $[-10L, 10L]$ and $M$ satisfy the conditions of Proposition \ref{Proposition: Trilinear Free Estimate}, with
    \[
    K := 
\begin{cases}
    |I_1|\vee M & \mbox{if }\eqref{Condition: M condition}\\
    |I_2| &\mbox{else}.
\end{cases}
    \]
    Then, for
    \[
    N_{13} = \mathrm{dist}(I_1, I_3)\quad\mbox{ and }\quad N_{23} = \mathrm{dist}(I_2, I_3)
    \]
    with $K\ll \max(N_{13}, N_{23})$, we have the following trilinear Strichartz estimate:
    \[
    \|P_{I_1}u_1P_{I_2}u_2P_{I_3}u_3\|_{L^2_{x,t}(\mu\mathbb{T}\times\lambda\mathbb{T})}\lesssim \log^3(\langle \lambda L\rangle)\left(\frac{1}{N}+\frac{K}{\max(N_{13}, N_{23})}\right)^\frac{1}{2}\prod_{j=1}^3\|P_{I_j}u_j\|_{Y^0_{+,\mu}}.
    \]
    Moreover, equivalent statements hold for conjugates so long as one has the analogous statement for $-I_j$.
\end{proposition}

\begin{corollary}\label{Corollary: Trilinear Strichartz Specific}
    Suppose that $s < \tfrac{\alpha}{2}$, $L\in 2^{\mathbb{N}_0}$, and $I_j\subset [-10L, 10L]$ for $j = 1, 2,3$ are intervals satisfying $|I_1|\lesssim |I_2|\lesssim |I_3|$. Suppose further that
    \[
    \mathrm{dist}(I_1, I_3)\sim \mathrm{dist}(I_2, I_3)\sim L,
    \]
    and that $|I_1|\ll L$. Then we have the following trilinear Strichartz estimate:
    \[
    \|P_{I_1}u_1P_{I_2}u_2P_{I_3}u_3\|_{L^2_{x,t}(\mu\mathbb{T}\times\lambda\mathbb{T})}\lesssim \log^3(\langle \lambda L\rangle)\left(\frac{1}{N}+\frac{|I_1|}{L}\right)^\frac{1}{2}\prod_{j=1}^3\|P_{I_j}u_j\|_{Y^0_{+,\mu}},
    \]
    with analogous results for conjugates so long as there are analogous statements with $-I_j$.
\end{corollary}

    \begin{equation}\label{Equation: Regularity}
        |\partial_{\ell_j}^\alpha \overline{M}_I(\ell_1,\dots,\ell_6)|\lesssim \frac{\cal{M}(I_1,\dots,I_6)}{L^\alpha_j}
    \end{equation}

\begin{lemma}\label{Lemma: M6 estimate}
    Let $I_1\dots, I_6$ denote intervals of length $L_i\in 2^\mathbb{Z}$. The multiplier 
    \[
    \overline{M}_6(k_1,\dots, k_6)\chi_{I_1}(k_i)\cdots\chi_{I_6}(k_6)
    \]
    defined on $\Gamma_6$ extends to $\mathbb{R}^6$ and satisfies size and regularity assumptions under the following conditions on $I_1,\dots I_6$.
    \begin{enumerate}[label = (\roman*)]
        \item For $|k_i|\sim N_i\in 2^{\mathbb{N}_0}$ and $N_1^*\sim N_2^*$ we have the size estimate and regularity estimates \eqref{Equation: Regularity} with 
        \begin{equation}\label{equation: M bound 1}
            |\cal{M}(I_1,\dots,I_6)|\lesssim m(N_1^*)N_1^*m(N_3^*)N_3^*.
        \end{equation}
        The same holds if $N_1^*\gg N_3^*$ and $|I_1|\sim |I_2|\sim N_3^*$.
        \item If $N_1^*\sim |k_1|\sim |k_2|\gtrsim N\gg N_3^*\sim N_4^*$, $|I_1|\sim |I_2|\lesssim \tfrac{(N_3^*)^2}{N_1^*}$, and $|k_1+k_2|\lesssim \tfrac{(N_3^*)^2}{N_1^*}$ for $k_j\in I_j$, then we have the size and regularity estimate with
        \begin{equation}\label{Equation: M bound 2}
            |\cal{M}(I_1,\dots, I_6)|\lesssim (N_3^*)^2.
        \end{equation}
        \item If $|I_j|\lesssim N_5^*$ for $j=1,2,3,4$, and $\max(|k_1+k_2|,|k_3+k_4|)\lesssim N_5^*$, then we have the size and regularity estimate with
        \begin{equation}\label{Equation: M bound 3}
            |\cal{M}(I_1,\dots, I_6)|\lesssim m(N_1^*)N_1^*m(N_5^*)N_5^*.
        \end{equation}
        \item If $|I_j|\lesssim N_{12}$ for $j = 1,2,3,4$, and $N_1\gtrsim N_{12}\sim |k_1+k_2|\sim |k_3+k_4|\gg |k_5^*|$, then the size and regularity estimate hold with
        \begin{equation}\label{Equation: M bound 4}
            |\cal{M}(I_1,\dots, I_6)\lesssim m(N_1^*)N_1^*m(N_{12})N_{12}.
        \end{equation}
    \end{enumerate}
\end{lemma}

\begin{remark}\label{Remark: N12 M term}
    The factor of $m(N_{12})$ in \eqref{Equation: M bound 4} is necessary to extending Theorem \ref{Theorem: GWP theorem} to $s < \tfrac{1}{3}$. In fact, this allows us to conclude that
    \[
    \frac{N^s}{N_1^s}\sum_{N_{12}\lesssim N_1}\frac{m(N_{12})N_{12}}{N} \lesssim 1 + \frac{N^s}{N_1^s}\sum_{N\ll N_{12}\lesssim N_1}N^{-s}N_{12}^s\lesssim 1,
    \]
    for some $N_1\gtrsim N$. This means we may utilize another factor of $N_1^{-s}$ to cancel this factor, and the above factor of $N^{-s}$ to cancel the $N^s$ factor that is now without a partner. 
\end{remark}

\begin{lemma}\label{Lemma: M6 integral term}
If $s > 0$ and $T\lesssim \mu = \lambda^{2-\alpha}N^{-\alpha}$, then
\begin{equation}\label{Equation: L6 lemma integral bound}
\left|\int_0^T\Lambda_6(\overline{M}_6)\,dt\right|\lesssim N^{-3+}\|Iu\|_{Y^1_T}^6.
\end{equation}
\end{lemma}

\begin{proof}
    We dyadically decompose into $|k_i|\sim N_i$, and let $N_j^*$ denote the decreasing rearrangement of the terms. By symmetry we may assume $N_1\geq N_2$ and $N_1\geq N_3\geq N_5$, $N_2\geq N_4\geq N_6$, and our assumptions on resonance afford us $\{N_1^*, N_2^*\} = \{N_1, N_2\}$. We again denote $\widehat{w}(\xi,t) = m(\xi)\langle\xi\rangle \widehat{u}(\xi,t)$, and reiterate that all even indexed terms implicitly come with complex conjugation. With this notation, the fundamental building block will then be the estimate
    \begin{multline}\label{Equation: Resonant Base Estimate}
        \left|\int_{\Gamma_6}\int_0^T\frac{\overline{M}_6(k_1,\dots, k_6)\prod_{j=1}^6\widehat{w}(k_j,t)}{\prod_{j=1}^6m(k_j)\langle k_j\rangle}\,dtd\Gamma_6\right|\\
        \lesssim \sum_{\substack{I_1,\dots, I_6\\I_i\subset A_{N_i}\\(\mathcal{R})\,\mathrm{ holds}}}\frac{\overline{\mathcal{M}}_6(I_1,\dots, I_6)}{\prod_{j=1}^jm(N_j)N_j}\left|\int_{\Gamma_6}\int_0^T\prod_{j=1}^6\widehat{w}(k_j,t)\chi_{I_j}(k_j)\,dtd\Gamma_6\right|,
    \end{multline}
    where we obtain bounds for $\overline{\mathcal{M}}_6$ via Lemma \ref{Lemma: M6 estimate}. This lemma will also provide us the regions $I_j$ on which we will sum. In most of the scenarios we have that the blocks $I_j$ are the dyadic $N_j$, but in $A_1$ and portions of $A_3$ they will be slightly smaller in accordance with Lemma \ref{Lemma: M6 estimate}.

    As is standard, we split the integral into regions
    \begin{align*}
        A_1 &= \{k\in\Gamma_6\,:\,N_2^*\gtrsim N\gg N_3^*\sim N_4^*\},\\
        A_2 &= \{k\in \Gamma_6\,:\, N_3^*\gtrsim N\gg N_4^*\},\\
        A_3 &= \{k\in\Gamma_6\,:\,N_4^*\gtrsim N\gg N_5^*\},\\
        A_4 &= \{k\in\Gamma_6\,:\, N_5^*\gtrsim N\},
    \end{align*}
    where there is no contribution from $A_2$ in this integral.\\
    
    \textbf{Case $A_1$: } In this region we take the $\overline{\mathcal{M}}_6$ bound of Equation \eqref{Equation: M bound 2}, $\overline{\mathcal{M}}_6\lesssim (N_3^*)^2$. We observe that we may apply the trilinear Strichartz of Corollary \ref{Corollary: Trilinear Strichartz Specific}. Hence, we find
    \begin{multline*}
        \sum_{\substack{I_1,\dots, I_6\\I_i\subset A_{N_i}\\(\mathcal{R})\,\mathrm{ holds}}}\frac{\overline{\mathcal{M}}_6(I_1,\dots, I_6)}{\prod_{j=1}^jm(N_j)N_j}\left|\int_{\lambda\mathbb{T}}\int_0^T\prod_{j=1}^6w_{N_j^*}\,dtdx\right|\\
        \lesssim \sum_{N_1^*\sim N_2^*\gg N_3^*\gtrsim\cdots\gtrsim N_6^*}\frac{N_3^*N_4^*}{N^{2-2s}(N_1^*)^{2s-}\prod_{j=3}^6N_j^*}\|w_{N_1^*}w_{N_3^*}w_{N_5^*}\|_{L^2_{x,t}}\|w_{N_2^*}w_{N_4^*}w_{N_6^*}\|_{L^2_{x,t}}\\
        \lesssim \sum_{N_1^*\sim N_2^*\gg N_3^*\gtrsim\cdots\gtrsim N_6^*}\frac{N_3^*N_4^*}{N^{2-2s}(N_1^*)^{2s-}N_3^*N_4^*(N_5^*N_6^*)^\frac{1}{2}}\left(\frac{1}{N_1^*} + \frac{1}{N}\right)\prod_{j=1}^6\|w_{N_j^*}\|_{Y^0}\\
        \lesssim \sum_{N_1^*\sim N_2^*\gg N_3^*\gtrsim\cdots\gtrsim N_6^*}N^{-2+}\left(\frac{1}{N_1^*} + \frac{1}{N}\right)\prod_{j=1}^6\|w_{N_j^*}\|_{Y^0}\lesssim N^{-3+}\|Iu\|_{Y^1_T}^6,
    \end{multline*}
    as desired. This also provides the blueprint for the remaining estimates, in that we will seek to always apply two trilinear estimates, unless we have enough high frequency terms to compensate.\\

    \textbf{Case $A_3$: } We further split $A_3$ into two regions, given as
    \begin{align*}
        A_{31} &= \{k\in A_3\,:\,N_1^*\gg N_3^*\}\\
        A_{32} &= \{k\in A_3\,:\, N_1^*\sim N_3^*\}.
    \end{align*}
    \textbf{Case $A_{31}$:} In this case we observe that
    \[
    |k_1^*\pm k_3^*|,\,|k_2^*\pm k_4^*|\gtrsim N_1^*,
    \]
    and that we may perform the almost orthogonal decomposition of $k_j\in I_j$ for $j = 1,2$ with $|I_j|\sim N_3^*$ so that we have the estimate $\overline{\mathcal{M}}_6\lesssim m(N_1^*)N_1^*m(N_3^*)N_3^*$ by equation \eqref{equation: M bound 1}. This will then follow from the observation that we may perform two applications of the trilinear Strichartz estimate of Corollary~\ref{Corollary: Trilinear Strichartz Specific}, which yields
    \begin{multline*}
    \|w_{N_5^*}w_{N_3^*}P_{I_1}w_{N_1^*}\|_{L^2_{x,t}}\|w_{N_6^*}w_{N_4^*}P_{I_2}w_{N_2^*}\|_{L^2_{x,t}}\\
    \lesssim (N_1^*)^{0+}\left(\frac{N_5^*}{N_1^*}+\frac{1}{N}\right)\prod_{j=1}^2 \|P_{I_j}w_{N_j^*}\|_{Y^0}\prod_{j=3}^6\|w_{N_j^*}\|_{Y^0},
    \end{multline*}
    and hence
    \begin{multline*}
        \eqref{Equation: Resonant Base Estimate}\lesssim \sum_{\substack{N_1^*\sim N_2^*\gg N_3^*>\dots> N_6^*\\I_1\sim I_2\\|I_1|,\,|I_2|\sim N_3^*}}\frac{N^{-2+2s}(N_{2}^*N_4^*)^{-s}}{N_5^*N_6^*}\left|\int_{\lambda\mathbb{T}}\int_0^T\prod_{j=1}^2P_{I_1}w_{N_j^*}\prod_{j=3}^6w_{N_j^*}\,dtd\Gamma_6\right|\\
        \lesssim \sum_{\substack{N_1^*\sim N_2^*\gg N_3^*>\dots> N_6^*\\I_1\sim I_2\\|I_1|,\,|I_2|\sim N_3^*}}\frac{N^{-2+2s}(N_{2}^*N_4^*)^{-s+}}{N_5^*N_6^*}\|w_{N_5^*}w_{N_3^*}P_{I_1}w_{N_1^*}\|_{L^2_{x,t}}\|w_{N_6^*}w_{N_4^*}P_{I_2}w_{N_2^*}\|_{L^2_{x,t}}\\
        \lesssim\sum_{N_1^*\sim N_2^*\gg N_3^*>\dots> N_6^*}\frac{N^{-2+2s}(N_{2}^*N_4^*)^{-s+}}{(N_5^*N_6^*)^{\frac12}}\left(\frac{1}{N_1^*} + \frac{1}{N}\right)\prod_{j=1}^6 \|Iu_{N_j^*}\|_{Y^1}\\
        \lesssim N^{-3+}\|Iu\|_{Y^1_T}^6.
    \end{multline*}
    
    \textbf{Case $A_{32}$:} Turning our attention to $A_{32}$ we have that $N_3^*\sim N_4^*$ and hence $N_1^*\sim N_4^*$, which makes the application of two trilinear estimates considerably more complicated.

    To overcome this, we decompose $A_{32}$ into
    \begin{align*}
        A_{321} &= \{k\in A_{32}\,:\,\{N_1^*,N_2^*,N_3^*,N_4^*\} = \{N_1,N_2,N_3,N_4\}\}\\
        A_{322} &= \{k\in A_{32}\,:\,\{N_1^*,N_2^*,N_3^*,N_4^*\} = \{N_1,N_2,N_3,N_5\}\}\\
        A_{321} &= \{k\in A_{32}\,:\,\{N_1^*,N_2^*,N_3^*,N_4^*\} = \{N_1,N_2,N_4,N_6\}\}.
    \end{align*}

    \textbf{Subcase $A_{321}$: } To utilize the trilinear Strichartz estimate here, we need to perform our (final) decomposition into the sets:
    \begin{align*}
        A_{3211} &= \{k\in A_{321}\,:\,k_1 >0,\, k_2 > 0,\, k_3 < 0,\, k_4 < 0\}\\
        A_{3212} &= \{k\in A_{321}\,:\,k_1 >0,\, k_2 < 0,\, k_3 < 0,\, k_4 > 0\}\\
        A_{3213} &= \{k\in A_{321}\,:\,k_1 >0,\, k_2 < 0,\, k_3 < 0,\, k_4 < 0\}\\
        A_{3214} &= \{k\in A_{321}\,:\,k_1 >0,\, k_2 < 0,\, k_3 > 0,\, k_4 < 0\},
    \end{align*}
    where we have used symmetry to assume that $k_1 > 0$. Note that we've possibly lost the original order of $N_1\geq N_3$ and $N_2\geq N_4$ in reducing the total number of possibilities to the above four, but the proof will be insensitive to this.

    We first observe that in either $A_{3211}$ and $A_{3212}$ we can perform two trilinear Strichartz estimates on the terms $u_1u_3u_5$ and $u_2u_4u_6$. Utilizing the bound from Equation \eqref{equation: M bound 1}, we have the size and regularity estimate
    \[
    \overline{\mathcal{M}}(I_1,\dots, I_6)\lesssim m(N_1^*)N_1^*m(N_3^*)N_3^*.
    \]
    Observing that
    \[
    |k_1-k_3|\sim |k_2-k_4|\sim |k_5-k_1|\sim |k_6-k_2|\sim N_1^*,
    \]
    we see that we may apply Corollary \ref{Corollary: Trilinear Strichartz Specific} twice, obtaining
    \begin{multline*}
        \eqref{Equation: Resonant Base Estimate}\lesssim \sum_{N_1\sim N_4\gtrsim N\gg N_5^*\geq N_6^*}\frac{N^{2s-2}(N_1)^{-2s}}{N_5N_6}\left|\int_{\lambda\mathbb{T}\times [0,T]}\prod_{j=1}^6w_{N_j^*}\,dxdt\right|\\
        \lesssim \sum_{N_1^*\sim N_4^*\gtrsim N\gg N_5^*}\frac{N^{2s-2}(N_1^*)^{-2s}}{N_5N_6}\|w_{N_1}w_{N_3}w_{N_5}\|_{L^2_{x,t}}\|w_{N_2}w_{N_4}w_{N_6}\|_{L^2_{x,t}}\\
        \lesssim \sum_{N_1\sim N_4\gtrsim N\gg N_5^*\geq N_6^*}\frac{N^{2s-3}(N_1^*)^{-2s+}}{(N_5N_6)^{\frac12}}\prod_{j=1}^6\|Iu_{N_j^*}\|_{Y^1}\lesssim N^{-3+}\|Iu\|_{Y^1_T}^6.
    \end{multline*}
    Suppose now that we are in case $A_{3213}$, in which case we have
    \[
    |k_2+k_3| = |k_1+k_4 + O(N_5^*)|\gtrsim N_1^*\gtrsim N,
    \]
    so that $|k_1+k_4|\gtrsim N_1^*$. In particular, we have access to two trilinear Strichartz applications and can apply the above argument again after a permutation of the functions.

    We are thus left with $A_{3214}$, where we distinguish two (further) cases:
    \begin{align*}
        A_{3214a} &:= \{k\in A_{3214}\,:\,|k_1+k_2|\lesssim |k_5^*|\}\\
        A_{3214b} &:= \{k\in A_{3214}\,:\,|k_1+k_2|\gg |k_5^*|\gtrsim |k_5+k_6|\}.
    \end{align*}
    In case $A_{3214a}$ we have that the linear relationship between $k_j$ forces $|k_3+k_4|\lesssim |k_5^*|$, and hence we may perform the further almost orthogonal decomposition of $k_j\in I_j$ for $j=1,2,3,4$ and $|I_j|\sim N_5^*$. Equation \eqref{Equation: M bound 3} yields the bound $\overline{\mathcal{M}}_6\lesssim m(N_1^*)N_1^*m(N_5^*)N_5^*$ in this scenario, and hence
    \begin{multline*}
        \eqref{Equation: Resonant Base Estimate}\lesssim \sum_{N_1\sim N_4\gtrsim N\gg N_5^*}\sum_{\substack{I_1\sim I_2\\ I_3\sim I_4\\|I_j|\sim N_5^*}}\frac{N^{3s-3}(N_1^*)^{-3s}}{N_6^*}\left|\int_{\lambda\mathbb{T}\times [0,T]}\prod_{j=1}^4P_{I_j}w_{N_j}\prod_{j=5}^6w_{N_j}\,dxdt\right|\\
        \lesssim \sum_{N_1\sim N_4\gtrsim N\gg N_5^*}\sum_{\substack{I_1\sim I_2\\ I_3\sim I_4\\|I_j|\sim N_5^*}}\frac{N^{3s-3}(N_1^*)^{-3s}}{N_6^*}\prod_{j=1}^4\|P_{I_j}w_{N_j}\|_{L^6}\prod_{j=5}^6\|w_{N_j}\|_{L^6}\\
        \lesssim N^{-3+}\|Iu\|_{Y^1_T}^6.
    \end{multline*}

    This leaves us with the final case of $A_{3214b}$, where $|k_1+k_2|\gg N_5^*\gtrsim |k_5+k_6|$. This forces $|k_3+k_4|\sim |k_1+k_2|$, so we perform a dyadic decomposition of $|k_1+k_2|\sim N_{12}$, and the further almost orthogonal decomposition of $k_j\in I_j$ for $|I_j|\sim N_{12}$ and $j = 1,2,3,4$. We note the bound $\overline{\mathcal{M}}_6\lesssim m(N_1^*)N_1^*m(N_{12})N_{12}$ of Equation \eqref{Equation: M bound 4}, which reduces us to bounding
    \begin{multline*}
        \eqref{Equation: Resonant Base Estimate}\lesssim \sum_{N_1^*\sim N_4^*\gtrsim N\gg N_5^*}\sum_{N_{12}\lesssim N_1^*}\sum_{\substack{I_1\sim I_2\\I_3\sim I_4\\|I_j|\sim N_{12}}}\frac{N^{3s-3}(N_1^*)^{-3s}m(N_{12})N_{12}}{N_5N_6}\\
        \times\left|\int_{\lambda\mathbb{T}\times [0,T]}\prod_{j=1}^4P_{I_j}w_{N_j}\prod_{j=5}^6w_{N_j^*}\,dxdt\right|.
    \end{multline*}
    In order to handle the $6$-linear estimate, we consider first the case that
    \[
    N_{12}^2\lesssim N_5^*N_1,
    \]
     and apply the base trilinear estimate of Proposition \ref{Lemma: trilinear Estimates} twice. This affords us
    \[
    \left|\int_{\lambda\mathbb{T}\times [0,T]}\prod_{j=1}^4P_{I_j}w_{N_j}\prod_{j=5}^6w_{N_j^*}\,dxdt\right|\lesssim (N_1^*)^{0+}\left(\frac{N_{12}}{N_1}+\frac{1}{N}\right)\prod_{j=1}^4\|P_{I_j}w_{N_j}\|_{Y^0}\prod_{j=5}^6\|w_{N_j}\|_{Y^0},
    \]
    where we may further sum to obtain
    \[
    \sum_{N_{12}^2\lesssim N_5^*N_1}(N_1^*)^{0+}m(N_{12})N_{12}\left(\frac{N_{12}}{N_1}+\frac{1}{N}\right)\lesssim (N_1^*)^{0+}N_5^* + N^{-s}N_1^{s+},
    \]
    by Remark \ref{Remark: N12 M term} and the summation restriction. If, on the other hand, we have
    \[
    N_5^*N_1\ll N_{12}^2,
    \]
    then we may apply the enhanced version of Proposition \ref{Lemma: trilinear Estimates} where
    \[
    K\lesssim \frac{N_5^*N_1}{N_{12}},
    \]
    and hence
    \[
    \left|\int_{\lambda\mathbb{T}\times [0,T]}\prod_{j=1}^4P_{I_j}w_{N_j}\prod_{j=5}^6w_{N_j^*}\,dxdt\right|\lesssim (N_1^*)^{0+}\left(\frac{N_5^*}{N_{12}}+\frac{1}{N}\right)\prod_{j=1}^4\|P_{I_j}w_{N_j}\|_{Y^0}\prod_{j=5}^6\|w_{N_j}\|_{Y^0}.
    \]
    Summation again yields the bound
    \[
    \sum_{N_5^*N_1\ll N_{12}^2\lesssim N_1^2}(N_1^*)^{0+}m(N_{12})N_{12}\left(\frac{N_{5}^*}{N_{12}}+\frac{1}{N}\right)\lesssim (N_1^*)^{0+}N_5^* + N^{-s}N_1^{s+}.
    \]
    It follows that in either situation we have
    \begin{multline*}
        \sum_{N_1^*\sim N_4^*\gtrsim N\gg N_5^*}\sum_{N_{12}\lesssim N_1^*}\sum_{\substack{I_1\sim I_2\\I_3\sim I_4\\|I_j|\sim N_{12}}}\frac{m(N_{12})N_{12}}{N^{3-3s}N_1^{3s-}N_5N_6}
        \left|\int_{\lambda\mathbb{T}\times [0,T]}\prod_{j=1}^4P_{I_j}w_{N_j}\prod_{j=5}^6w_{N_j^*}\,dxdt\right|\\
        \lesssim \sum_{N_1^*\sim N_4^*\gtrsim N\gg N_5^*}N^{-3+}N_1^{0-}\prod_{j=1}^6 \|w_{N_j}\|_{Y^0}\lesssim N^{-3+}\|Iu\|_{Y^1}^6,
    \end{multline*}
    which is our desired bound.

    \textbf{Subcases $A_{322}$ \& $A_{323}$: }We assume by symmetry of the conditions of resonance that we are in case $A_{322}$. We note that we must have that $k_1,k_3,k_5$ are not of the same sign. Moreover, if $k_i$ and $k_2$ are of the same sign then $|k_i-k_2|\sim N_1^*$, and $|k_i+k_2|\sim N_{1}^*$ if they have opposite signs. 
    
    We can assume without loss of generality that $k_1 > 0$. We then have the following cases:
    \begin{enumerate}[label = (\roman*)]
        \item $k_3 > 0,\,k_5 < 0$ or $k_3 < 0,\,k_5 > 0$
        \item $k_3 < 0$ \& $k_5 < 0$.
    \end{enumerate}
    We first note that regardless of the sign of $k_2$, we always have $|k_j+k_2|\gtrsim N_1^*$ for $j\in\{1,3,5\}$. Indeed, if they both have the same sign then it's immediate, while if they differ in signs the resonance condition will force $|k_j+k_2|\gtrsim N_1^*$. Now, if we fall in case $(i)$ then we have
    \[
    |k_1+k_2|,\,|k_3-k_5|\gtrsim N_1^*,
    \]
    so that we may apply two trilinear estimates to $w_{N_1}w_{N_2}w_{N_4}$ and $w_{N_3}w_{N_5}w_{N_6}$. If we fall in case $(ii)$ then
    \[
    |k_1-k_3|,\,|k_2+k_5|\gtrsim N_1^*,
    \]
    and we may apply two trilinear estimates to $w_{N_1}w_{N_3}w_{N_4}$ and $w_{N_2}w_{N_5}w_{N_6}$.

    Thus, in either case we may apply the bound $\overline{\mathcal{M}}_6\lesssim m(N_1^*)N_1^*m(N_3^*)N_3^*$ together with the two applications of Corollary \ref{Corollary: Trilinear Strichartz Specific} to obtain the bound
    \[
\eqref{Equation: Resonant Base Estimate}\lesssim N^{-3+}\|Iu\|_{Y^1_T}^6,
    \]
    as in the proof of $A_{3211}$.\\
    \textbf{Case $A_4$:} This is the final case, in which all frequencies are large. Here we just assume that $N_1\sim N_2\geq N_3\geq \dots\geq N_6$ and take the $\overline{\mathcal{M}}_6$ bound afforded by \eqref{equation: M bound 1}:
    \[
    \overline{\mathcal{M}}\lesssim m(N_1)N_1m(N_3)N_3.
    \]
    In this manner, we find via a straightforward $L^6$ estimate:
    \[
        \eqref{Equation: Resonant Base Estimate}\lesssim \sum_{N_1\sim N_2\geq N_3\dots N_6}N^{-3+3s}(N_2N_4N_5)^{-s}\prod_{j=1}^6\|w_{N_j}\|_{L^6_{x,t}}\lesssim N^{-3+}\|Iu\|_{Y^1_T}^6,
    \]
    which was our desired estimate.
    % Here Lemma \ref{Lemma: M6 estimate} implies that we may perform the almost orthogonal decomposition $k_j\in I_j$ with $|I_j| = N_3^*$ for $j = 1,2$ and also that $\overline{\mathcal{M}}_6\lesssim (N_3^*)^2$. Thus two applications of the bilinear Strichartz estimate concludes the proof
\end{proof}

\end{document}
